\documentclass[10pt,a4paper]{amsart}
\usepackage[margin=19mm]{geometry}
\usepackage{amsmath,amssymb,mathtools}
\usepackage[scr=rsfs,cal=euler]{mathalfa}
\usepackage{xfrac}
\usepackage[unicode,hidelinks,bookmarks=false]{hyperref}
\newcommand{\ie}{i.e.\ }
\newcommand{\cf}{cf.\ }
\newcommand{\resp}{respectively\ }
\newcommand{\Subsection}{\subsection}
\providecommand{\nobreakdash}{\nobreak\hbox{-}}
\setlength{\emergencystretch}{2em}
\multlinegap0pt
\usepackage{bm}
\usepackage{bbm}
\usepackage{dsfont}
\usepackage{xspace}
\usepackage{cancel}
\usepackage{mathtools}
\usepackage{amsthm}
\makeatletter
\@ifundefined{corollary}{\newtheorem{corollary}{Corollary}}{}
\@ifundefined{lemma}{\newtheorem{lemma}[corollary]{Lemma}}{}
\makeatother
\newcommand{\R}{\mathbf{R}}
\title[Normal invariant of nearby Lagrangians via twisted derivativ]{Normal invariant of nearby Lagrangians via twisted derivative}
\author{Mohammed Abouzaid, Daniel \'Alvarez-Gavela, Sylvain Courte, Thomas Kragh}
\date{Source: arXiv:2505.12515v2 (CC BY 4.0)}
\begin{document}
\maketitle

\noindent\emph{Source and licence.} Normal invariant of nearby Lagrangians via twisted derivative. Version arXiv:2505.12515v2 (CC BY 4.0), distributed under the Creative Commons Attribution 4.0 International licence, as recorded in the arXiv licence metadata for this exact version and shown on the abstract page https://arxiv.org/abs/2505.12515v2. Full rights notice: licenses/arxiv-2505.12515-CC-BY-4.0.txt.\par\smallskip

\noindent\textit{Editorial scope.} Counted unit: the Lemma whose source label is \texttt{lem:propertyFa} in the article (source0.tex lines 937--946, source label \texttt{lem:propertyFa}), delivered together with the article's own complete original proof of it (source0.tex lines 947--987, one \texttt{proof} environment of 41 lines). No mathematics is written, repaired or re-derived: statement and proof are the article's own text line for line. Required context retained uncounted: none. Every definition, display and label of those lines is retained unchanged. Omitted from this package and not redistributed: 1--936, 988--2547. Nothing inside a retained excerpt is omitted or altered: every retained body is delivered exactly as the article's lines are written, so no omission receipt of any kind accompanies this package. Citations: the retained excerpts carry no citation command at all, so no bibliography entry of the article is transcribed and no reference is redistributed. Reference adaptation: none was needed and none was made; every reference inside the delivered unit resolves inside it, so no unresolved reference or citation is produced. Printed slips kept literal: no printed word, symbol, hypothesis or number of the counted statement or of its proof was altered. Layout and class: the article's own class is \texttt{amsart} with packages including amsrefs, amssymb, comment, etex, xy, mathtools, euscript, xcolor, tikz, pgfplots, tikz-cd; the wrapper is the lane's pinned \texttt{amsart} editorial wrapper at 10pt on A4, which restates the theorem-like environments the retained text needs and adds the article's own macro definitions that the retained text needs (\texttt{\textbackslash R}), with extra wrapper packages (bm, bbm, dsfont, xspace, cancel, mathtools). The delivered numbering is the unit's own. Assets: no retained span contains a figure environment, a graphics-inclusion command, a PDF-page inclusion, a table or a TikZ picture; no image file is copied into this package, the package declares no asset dependency and no image file is redistributed with it. Licence: version arXiv:2505.12515v2 of this article is distributed under the Creative Commons Attribution 4.0 International licence, as recorded in the arXiv licence metadata for this exact version and shown on the abstract page https://arxiv.org/abs/2505.12515v2. A full rights notice is delivered at licenses/arxiv-2505.12515-CC-BY-4.0.txt. Characters: every retained excerpt is pure ASCII, so the delivered engine, XeLaTeX with its default Latin Modern text font, reports no missing character.
\begin{lemma}\label{lem:propertyFa}
% The function $F_a$ satisfies the following properties:
% \begin{enumerate}
% \item For all $x\in M$, $F_a^x$ is homogeneous of degree $2$, namely for all $\lambda \geq 0$
% and $(t,w,v)\in \R^{2+n}$, $F_a^x(\lambda t,\lambda w, \lambda v)=\lambda^2 F^x_a(t,w,v)$,
% \item For all $a\geq c$, $F_a$ is $C^1$ on $M'\times \R^{2+n}$,
% \item
  The restriction of $G^x_a$ to $\R^{2+n}\setminus\{0\}$ is transverse to $0$ if and only if the zero-level set of $f^x$ is transverse to the hypersurface $ \frac{w^2}{a^2}+\frac{v^2}{c^2} = 1$.
%\end{enumerate}
\end{lemma}

\begin{proof}
%\begin{enumerate}
% \item Homogeneity is clear from the formula for $F_a$.
% \item On the unit sphere sphere $\{r=1\}$ we have
% \[F^x_a(t,w,v)=\frac{1}{a}\left(aw+\epsilon^x(aw,cv)\delta_{t>0}+g^x(cv)\right).\]
% If $a\geq c$, we have $|(aw,cv)|\geq c\sqrt{1-t^2}$. This shows, according to our assumption on $c$,
% that $(aw,cv)$ leaves the support of $\epsilon$
% as $t$ goes to $0^+$, and this is valid for all $x\in K$. So $F_a$ is smooth on $M'\times S^{n-1}$
% and therefore $C^1$ on $M'\times \R^n$ according to Lemma~\ref{lem:radialC1}.
% \item
 % We study $F^x_a$ separately on two open sets which cover $\R^{2+n}\setminus \{0\}$.
We first prove transversality in the regions $\{0 < t\}$ and $\{t > 0\} $, which can both be equipped with coordinates 
\begin{equation} \label{eq:change_coordinates}
    \left(s=r^2,w'=\frac{aw}{r},v'=\frac{cv}{r}\right)
\end{equation}
which define a diffeomorphism between these regions and the product of $(0,+\infty)$ with the interior of the ellipsoid
\begin{equation}
    (0,+\infty)\times \left\{\frac{w'^2}{a^2}+\frac{v'^2}{c^2} < 1\right\}.
\end{equation}
%$\{t>0\}\to $.
% with inverse 
% \[\left(t=\sqrt{s}\sqrt{1-\frac{(w')^2}{a^2}-\frac{|v'|^2}{c^2}},w=\sqrt{s}\frac{w'}{a},v=\sqrt{s}\frac{v'}{c}\right).\]
In these coordinates, the function $G^x_a$ reads simply:
\begin{equation} \label{eq:expressing_F_in_half_space}
  G^x_a(s,w',v') =
  \begin{cases}
    \frac{s}{a}f^x(w',v') &   \textrm{ for } 0 \leq t \\
    \frac{s}{a} \left( w' + g^x(v') \right) &  \textrm{ for } t \leq 0 .
  \end{cases}
\end{equation}
The first expression vanishes transversely since $f^x$ does, and the second because the $w'$ derivative is non-trivial.




It thus remains to ensure transversality along the hypersurface $t=0$, and we can use the closure of either chart for this argument:  in the region $0 \leq t$ the function $G^x_a$ is obtained by pulling back the restriction of $ \frac{s}{a}f^x(w',v')$ to the product
\begin{equation}
    (0,+\infty)\times \left\{\frac{w'^2}{a^2}+\frac{v'^2}{c^2} \leq  1\right\}
\end{equation}
under the smooth map given by Equation \eqref{eq:change_coordinates}. On the boundary, the $t$-derivative of this map identically vanishes, so that the partial derivative $\partial G^x_a / \partial t$ is trivial. This implies that critical points of $G^x_a$ which lie in the hypersurface $t=0$ are exactly the critical points of the restriction of this function. This hypersurface is diffeomorphically identified by Equation \eqref{eq:change_coordinates} with the product of $(0,\infty)$ with the hypesurface $ \frac{w^2}{a^2}+\frac{v^2}{c^2} = 1$, and the expression for $G^x_a$ in Equation \eqref{eq:expressing_F_in_half_space} again identifies the critical points of $G^x_a$ with the product with $(0,\infty) $ of the restriction of the critical points of $f^x(w',v')$ to the boundary of the ellipsoid.
\end{proof}

\end{document}
